
\chapter{User experiments}
This chapter presents the experiments conducted to evaluate the 2 contributions of the Miwo platform. The main focus is the evaluation of the expressive motion framework, but it also presents progress in RL locomotion. 

Evaluating the overall expressiveness of a robot is different from evaluating other robotic systems. The performance cannot be evaluated by objective metrics alone. If a motion is technically smooth and precise but fails to communicate the intended emotion or intent, it must be considered unsuccessful. To accurately evaluate the expressive capabilities of Miwo, we rely on a Human perception study in which participants are asked to classify and evaluate different expressive motions either by interacting directly with the physical robot or by watching short video clips of the robot performing each motion. 

In addition to the human perception trials, this chapter will also cover the experiments and results related to the RL locomotion. 


\section{Experiment design and Participants}
The experiment was designed to evaluate the MiWo robot's overall expressive capabilities through a human perception study. Participants initially encountered the robot in an idle state before being presented with 6 different expressive motions. After each animation, the participants were asked to classify what emotion or intent they believed the robot was expressing. Participants were also asked to evaluate the overall impression of the robot after all 6 motions were completed. This study employed a questionnaire-based design to collect responses, and quantitative methods were used to assess the clarity of expressive motions.

\subsection{Experiment setup}               
The participants of the study can be divided into two groups based on how they experienced the robot's expressive motions: 
The first group attended live demonstrations, where they observed the physical robot performing directly in front of them, allowing them to get a better impression of its size, mechanics, and weight. 
The second group completed the experiments remotely by watching short video clips of the same motions performed in the live demonstrations. 
By offering both in-person and remote participant options, the study was accessible to a broader range of participants. The responses of the 2 groups were also analyzed separately to identify any differences in the perception between live and video-based observations.
                
\subsection{Participants}
Participants were recruited using a wide variety of methods. The participants for the live demonstrations were mostly recruited by asking fellow students and researchers available on campus. Recruitment for the video-based study involved creating multiple social media posts, posting posters around campus, and sending direct invitations to the study to various friends and family, resulting in a wide range of participants with diverse backgrounds.
A total of 38 participants completed the study: 16 attended the live demonstration, and 22 completed the remote video-based study. The participants were not required to have any prior knowledge of robotics or engineering, and no demographic restrictions were imposed other than a minimum age requirement. Before participating in the study, all participants were informed of its purpose and provided consent. All responses were collected anonymously using a participant ID system that could not be linked back to individual participants, and no personal data was collected. Data was stored securely and used solely for academic research purposes in accordance with the institution's data management policy.

\begin{figure}[H]
\centering

\begin{subfigure}[t]{0.48\textwidth}
\centering
{\small\textbf{Live study (N=16)}}
\vspace{1.0cm}

\begin{tikzpicture}[scale=0.7]
\pie[
    color={navyblue, emotion},
    text=pin,
    before number=,
    after number=\%,
    font=\tiny,
    radius=2,
]
{88/Engineering, 12/Natural Sciences}
\end{tikzpicture}

\vspace{1.0cm}
{\small \textbf{Field of study}}
\vspace{0.8cm}

\begin{tikzpicture}[scale=0.7]
\pie[
    color={navyblue, emotion},
    text=pin,
    before number=,
    after number=\%,
    font=\tiny,
    radius=2,
]
{50/High familiarity, 50/Low familiarity}
\end{tikzpicture}

\vspace{0.5cm}
{\small \textbf{Robot familiarity}}
\end{subfigure}
\hfill
\begin{subfigure}[t]{0.48\textwidth}
\centering
{\small\textbf{Video study (N=22)}}
\vspace{0.1cm}

\begin{tikzpicture}[scale=0.7]
\pie[
    color={navyblue, emotion, adjacent, lightgrey, intended},
    text=pin,
    before number=,
    after number=\%,
    font=\tiny,
    radius=2,
]
{50/Engineering, 14/Social Sciences,
 9/Medicine, 9/Natural Sciences, 18/Other}
\end{tikzpicture}

\vspace{0.1cm}
{\small \textbf{Field of study}}
\vspace{0.8cm}

\begin{tikzpicture}[scale=0.7]
\pie[
    color={navyblue, emotion},
    text=pin,
    before number=,
    after number=\%,
    font=\tiny,
    radius=2,
]
{41/High familiarity, 59/Low familiarity}
\end{tikzpicture}

\vspace{0.6cm}
{\small \textbf{Robot familiarity}}
\end{subfigure}

\caption[Demographic breakdown of participants]{Demographic breakdown of participants. The top row shows
Field of study distribution, the bottom row shows robot familiarity
grouped into high (moderately familiar, very familiar, or expert)
and low (not familiar or slightly familiar).}
\label{fig:demographics_pie}
\end{figure}

% bar/pie chart for participants and demographics
As shown in \autoref{fig:demographics_pie}, the two groups differ noticeably in their respective backgrounds. The live study group, consisting of participants found on campus, is almost exclusively from engineering and technology fields, making up 88\% of participants. This displays potential weakness in the campus-based recruitment process. The video-based study shows a more diverse distribution, with only 50\% of participants doing engineering. Other differences include the live-study group reporting a higher level of overall robot familiarity, with 50\% rating themselves as having a higher level of familiarity compared to the 41\% in the video-based studies. These demographic differences are acknowledged as potential factors in interpreting the performance differences between the two studies and are addressed further in the results chapter.


\subsection{Procedure and questionnaire} 
In the first part, where participants experienced the robot live, they entered the room to find it already turned on and in an idle state, performing breathing-like, subtle movements. To provide a similar experience to the group in the video-based study, participants were first asked to watch a one-minute video of the robot standing and moving in its idle state, allowing them to familiarize themselves with the robot's idle behavior and appearance, ensuring that the following animations were evaluated relative to the robot's idle state rather than isolated movements without context. Each participant was then presented with 6 video clips or live demonstrations of the robot performing distinct expressive motions. The motions presented were designed to correspond to the following emotions or intents: Curiosity, Happiness, Sadness, Disappointment, Anger, and Fear. For each motion, participants were asked to answer 4 questions. The first question asked participants to name the emotion or intent the robot was expressing in their own words, with no predefined options, allowing for more open-ended responses. This was followed by 3 questions asking the participants to rate their confidence in their answers, how well they felt the robot expressed the named emotion, and how natural the motion felt on a scale of 1 to 5.

\begin{table}[H]
\centering
\begin{tabular}{|l|l|l|}
    \hline
    \textbf{Question} & \textbf{Type} & \textbf{Scale} \\
    \hline
    Q1. What emotion or intent does the robot & Open-ended & Free text \\
    appear to be expressing? & & \\
    \hline
    Q2. How confident are you in your answer? & Likert scale & 1 -- 5 \\
    \hline
    Q3. How well do you think the robot & Likert scale & 1 -- 5 \\
    expressed this emotion? & & \\
    \hline
    Q4. How natural did the robot's & Likert scale & 1 -- 5 \\
    movement feel? & & \\
    \hline
\end{tabular}
\caption{Questions the participants were faced with during the emotion/intent recognition part of the study.}
\label{tab:questionnaire}
\end{table}

The open-ended format for the question regarding emotion/intent recognition was deliberately chosen to avoid influencing participants into a predefined set of options. By doing so, we allow for a more genuine assessment of whether the intended emotion/intent was communicated clearly without guiding the participant to a specific answer.

Following the individual evaluation, participants would provide an overall evaluation of the robot using the Godspeed scale. The Godspeed scale was introduced in \cite{bartneck_measurement_2008} as a standardized instrument for measuring human-robot interactions and is widely used and recognized. The Godspeed scale evaluates perceptions of the robot across 5 key concepts in HRI: anthropomorphism, animacy, likability, perceived intelligence, and perceived safety. Each concept is assessed using 5 consistent questionnaires that employ semantic differential scales. Using a standardized, validated metric for the overall evaluation allows the study's results to be compared with other work in expressive robotics and human-robot interaction.

The order of the motions was the same for each participant, and the questionnaire was administered digitally using Google Forms for both the live and video-based studies, taking approximately 10-15 minutes to complete.

\subsection{Metrics and evaluation}
It is worth noting that the study collects two categories of data. The emotion/intent recognition data consist of open-ended labels, which were analyzed by grouping responses into emotion/intent categories and calculating recognition rates for each animation. The scaled responses for analyzing confidence, expressiveness, and naturalness were analyzed quantitatively, while the godspeed scale data were scored and analyzed per dimension.

Due to the open-ended nature of the emotion/intent recognition question, responses were grouped into three categories for better visualization and quantitative analysis. The response groups include: intended emotion (response matching the intended emotion), adjacent emotion (response describing a related or overlapping emotional state), and other (responses unrelated to the intended emotion). Grouping participants' responses was necessary to enable quantitative comparisons across animations while preserving the insights of the open-ended response format. 
An example of the grouping logic can be seen in the curiosity animation. One participant classified the animation as \textbf{'cautious-}, this was classified as adjacent since caution shares the attentive qualities of the curiosity animation, even if it does not directly match the name of the intended emotion. 



